% First-reader proposal, scan 008; root adjudication pending.
\scanpage{8}\mspage{7}
Je voudrais maintenant faire passer les axiomes de disjonction-compatibilité avant ceux sur les lieux, pour ne pas \inserted{avoir l'air de} les en faire dépendre. Les axiomes sont ``facultatifs'' pour un atelier, et il y aura pas mal \unsure{des'|des|d'} implications mutuelles avec les axiomes sur les lieux.

Une notion essentielle ici est celle de ``$F$ et $G$ \uline{disjoints}''
\[
 F\parallel G\ \overset{\unsure{\mathrm{dfn}|\mathrm{defn}}}{\Longleftrightarrow}\ F\mathrel{\tikz[baseline=-.4ex,x=.65em,y=.65em]{\draw (0,1.1)--(1,.5)--(0,-.1) (1,1.1)--(0,.5)--(1,-.1);}}G\quad\text{et}\quad F\cap G=\varnothing_{\unsure{\mathscr F|\mathcal F|F}}
\]
notion qui ne dépend que de $\leq$. On a d'ailleurs
\[
 F\parallel G\ \Longleftrightarrow\ \forall X\in\widetilde F,\ Y\in\widetilde G,\quad X\parallel Y
\]
donc interprétation en \unsure{termes|terme} de $X\parallel Y$ dans $\unsure{M|\mathcal M}$, qui y signifie
\[
\begin{aligned}
X\mathrel{\underset{M}{\parallel}}Y\ \Longleftrightarrow\
 X\mathrel{\underset{M}{\mathrel{\tikz[baseline=-.4ex,x=.65em,y=.65em]{\draw (0,1.1)--(1,.5)--(0,-.1) (1,1.1)--(0,.5)--(1,-.1);}}}}Y\quad &\text{et}\quad\{X,Y\}\ \text{n'est \uline{pas}}\\
&\text{minoré dans } M,\leq\\
&\text{i.e. }\nexists Z\ \text{avec }Z\leq X,Z\leq Y
\end{aligned}
\]

(\uline{NB} Il \deleted{\unsure{résultera}}\unsure{résulte|résultera} des axiomes \inserted{précédents} que ça implique que $\nexists Z$ avec $Z\ll X,Z\ll Y$, \deleted{\unsure{i.e.} c'est un $\unsure{\exists}$} et \deleted{\unsure{la}} réciproque résultera \unsure{d'axiome|d'axiomes} ci-dessous\ldots)

% slanted
\margin{Un atelier sera dit \uline{spacieux} s'il satisfait cette condition. Contre-ex. : les ateliers simpliciaux\ldots}%
\deleted{\uline{At DC} 1 \quad (\uline{Axiome des ateliers spacieux}) \unsure{Soient} $X,Y\in M$, \inserted{\unsure{soit}} \deleted{\unsure{et}} $K=X\cap Y=\underset{Z\in\widetilde X\cap\widetilde Y}{\mathrm{Sup}}\,Z$. \deleted{\unsure{Il}} Supposons que \inserted{pour $X',Y'\in M$}}
\[
\deleted{\deleted{\unsure{\forall}}\ X'\ll X,\ X'\parallel K,\ Y'\ll Y,\ Y'\parallel K\ \Longrightarrow\ X'\parallel Y'.}
\]
\deleted{Alors [NB et alors seulement\ldots\ \inserted{cf Cor}] $X\mathrel{\tikz[baseline=-.4ex,x=.65em,y=.65em]{\draw (0,1.1)--(1,.5)--(0,-.1) (1,1.1)--(0,.5)--(1,-.1);}}Y$.}

\deleted{\uline{Cor} Soient $F,G\in\mathscr F$, $K=F\cap G=\underset{Z\in\deleted{\unsure{\widetilde X|X}}\unsure{\widetilde F|F}\cap\deleted{\unsure{\widetilde Y|Y}}\widetilde G}{\mathrm{Sup}}\,Z$.}
\[
\begin{aligned}
&\deleted{\text{Alors } F\mathrel{\tikz[baseline=-.4ex,x=.65em,y=.65em]{\draw (0,1.1)--(1,.5)--(0,-.1) (1,1.1)--(0,.5)--(1,-.1);}}G\quad \text{\tikz[baseline=(s.base)]{\node[inner sep=1pt] (s) {ses};\draw (s.center) ++(30:1.05em and .72em) arc[start angle=30,end angle=-295,x radius=1.05em,y radius=.72em];}}\quad \forall X,\deleted{\unsure{\varphi|y}}\unsure{Y|y}\in M}\\
&\deleted{X\ll F,\ Y\ll G,\ X\parallel K,\ Y\parallel K\ \Longrightarrow\ X\parallel Y}
\end{aligned}
\]
